\chapter{Hydration: Attaching HRBR to Server HTML}

\section*{What you'll build}

In this chapter I’ll build a mental model for v1 hydration in Relax, and you’ll learn how to reason about it by reading the real runtime code.
Concretely, you’ll be able to:

\begin{itemize}
  \item take an HTML string that was produced from a BlockDef, insert it into the DOM, and \emph{attach} a live HRBR block to it without rewriting nodes
  \item explain why Relax’s v1 hydration is strict and why that strictness is a feature (not a limitation)
\end{itemize}

\section*{What you'll learn}

\begin{itemize}
  \item The hydration contract: what must match between server HTML and \texttt{templateHTML}
  \item The two safety rails: \texttt{root tag check} and \texttt{slot-path validation}
  \item The cost model: why hydration is cheap when it succeeds, and cheap to bail out when it must
  \item How v1 hydration handles DOM updates (same slot kinds as Chapter~14) and where it differs (events)
\end{itemize}

\section*{Prereqs}

\begin{itemize}
  \item Chapter~14: Blocks, templates, and slots (slot kinds + path addressing)
  \item Chapter~12: Effects and dependency tracking (why \texttt{mountCompiledBlock} works)
\end{itemize}

\section{Why hydration exists (from scratch)}

Hydration is the bridge between two worlds:

\begin{itemize}
  \item The \textbf{server} produced HTML as a string. That’s great for first paint and SEO.
  \item The \textbf{client} needs live behavior: event handlers, reactive updates, and cleanup.
\end{itemize}

If you don’t hydrate, you have only two choices on the client:

\begin{itemize}
  \item \textbf{Remount by recreating DOM} from scratch, which throws away the server’s nodes (and can introduce a visible flash), or
  \item \textbf{Try to patch your way into the server DOM} with a general reconciler, which is powerful but hard to make airtight.
\end{itemize}

Relax HRBR v1 chooses a third option:

\begin{quote}
	extbf{Attach by validating structure and resolving slot node references.} If validation fails, remount.
\end{quote}

This conservative posture is intentional. It makes hydration predictable: the moment we say “hydrated”, we’re also saying the block’s template ABI matches the DOM we’re about to mutate.

\section{The big picture: where hydration sits in HRBR}

A block runtime has a natural relationship with server rendering:

\begin{itemize}
  \item The compiler produces a stable template shape (\texttt{templateHTML}).
  \item The server can render the template to HTML.
  \item The client can \emph{attach} to the existing DOM instead of rebuilding it.
\end{itemize}

In Relax HRBR v1, hydration is deliberately conservative. If anything looks wrong, it bails out and remounts client-side.

\begin{center}
\begin{tikzpicture}[
  node distance=9mm,
  box/.style={draw, rounded corners, align=left, inner sep=5pt},
  arrow/.style={-Latex, thick}
]
\node[box] (start) {\textbf{hydrateBlock(def, host)}\\try attach to server DOM};
\node[box, below=of start] (hasroot) {\textbf{host.firstElementChild?}\\no \textrightarrow mount fresh};
\node[box, below=of hasroot] (rootok) {\textbf{root tag matches?}\\no \textrightarrow clear + mount};
\node[box, below=of rootok] (slots) {\textbf{resolve + validate slots}\\(paths, node kinds, tagNames)};
\node[box, below=of slots] (ok) {\textbf{success}\\return hydrated block\\(identity preserved)};

\draw[arrow] (start) -- (hasroot);
\draw[arrow] (hasroot) -- (rootok);
\draw[arrow] (rootok) -- (slots);
\draw[arrow] (slots) -- (ok);
\end{tikzpicture}
\end{center}

This chapter is grounded in:

\begin{itemize}
  \item \texttt{runtime/hydration.ts} (hydrate implementation)
  \item \texttt{runtime/ssr.ts} (string rendering helper used in tests)
  \item \texttt{runtime/\_\_tests\_\_/hydration.test.ts}
  \item \texttt{runtime/\_\_tests\_\_/ssr.test.ts}
\end{itemize}

\section{A tiny contract (what hydrateBlock promises)}

Hydration is the moment where two worlds meet:

\begin{itemize}
  \item The \textbf{server} produced HTML as a string.
  \item The \textbf{client} wants a live block instance with slot node references.
\end{itemize}

In v1, Relax uses a strict approach that is easy to reason about and hard to get subtly wrong.
You can treat \texttt{hydrateBlock(def, host, initialValues)} as enforcing this contract:

\begin{enumerate}
  \item If there is no DOM to attach to (\texttt{host.firstElementChild == null}), hydration simply mounts a fresh block.
  \item If the root element tag does not match the template’s root tag, hydration clears the host and mounts fresh.
  \item If slot paths resolve to the wrong node types or to the wrong element tag names (compared to the parsed template), hydration clears the host and mounts fresh.
  \item If everything matches, hydration returns a block-like object whose \texttt{root} and \texttt{slotNodes} refer to the server DOM nodes.
\end{enumerate}

The headline design choice is: \textbf{when in doubt, remount}. That makes correctness obvious at the cost of being less clever than incremental/partial hydration systems.

\section{Build it: a minimal end-to-end hydration loop}

I like to teach hydration as a 3-step loop: \textbf{render} \(\rightarrow\) \textbf{attach} \(\rightarrow\) \textbf{update}.

In Relax, the pieces are:

\begin{itemize}
  \item \texttt{renderBlockToString(def, \{ initialValues \})} (SSR helper)
  \item \texttt{host.innerHTML = html}
  \item \texttt{hydrateBlock(def, host, initialValues)} (attach)
  \item optionally: \texttt{mountCompiledBlock(def, host, compiledSlots)} to wire signals/effects (Chapter~14)
\end{itemize}

Here’s the smallest "looks runnable" snippet that matches what the tests demonstrate.

\begin{lstlisting}[language=JavaScript, caption={SSR + hydration + update: the smallest end-to-end loop}]
import { defineBlock, mountCompiledBlock } from './runtime/block'
import { hydrateBlock } from './runtime/hydration'
import { createSignal } from './runtime/signals'
import { renderBlockToString } from './runtime/ssr'

const Counter = defineBlock({
  templateHTML: `<div class="c"><button id="inc">+</button><span id="t"> </span></div>`,
  slots: { value: { kind: 'text', path: [1, 0] } },
})

// "Server" side
const html = renderBlockToString(Counter, { initialValues: { value: 'A' } })

// "Client" side
const host = document.getElementById('host')
host.innerHTML = html

// Attach to existing nodes (no rewrite if markup matches)
const hydrated = hydrateBlock(Counter, host)

// Wire reactive updates on top (optional, but this is the HRBR story)
const [value, setValue] = createSignal('A')
const mounted = mountCompiledBlock(Counter, host, [{ key: 'value', read: () => value() }])

setValue('B')
mounted.update({ value: 'B' })

mounted.dispose()
hydrated.destroy()
\end{lstlisting}

\subsection{Why these steps matter}

\begin{itemize}
  \item \textbf{SSR produces a shape:} If SSR and the template disagree, hydration must not "half attach".
  \item \textbf{Hydration creates references:} the essential output is \texttt{slotNodes[key] = Node}.
  \item \textbf{Updates are the same contract:} a call to \texttt{update(\{ key: value \})} still means “write into slot nodes”.
\end{itemize}

\section{Hydration contract (v1)}

\texttt{hydrateBlock(def, host, initialValues?)} returns a \texttt{HydratedBlock} (type alias of \texttt{MountedBlock}).

Its docstring states three key assumptions:

\begin{enumerate}
  \item \texttt{host} contains exactly one root element matching the template.
  \item Hydration does not mutate DOM; it only resolves slot node references.
  \item On mismatch, hydration bails out and remounts client-side.
\end{enumerate}

This is a strict, "safety first" strategy: it prefers correctness over partial hydration tricks.

It’s also worth noticing what the contract does \emph{not} promise:

\begin{itemize}
  \item It does not try to recover from mismatches by patching individual nodes.
  \item It does not attempt to reconcile text/comment whitespace differences.
  \item It does not do event delegation or partial listener attachment.
\end{itemize}

\section{The fast path: attach without rewriting}

Given server markup like:

\begin{verbatim}
<div><span>A</span></div>
\end{verbatim}

and a matching block definition:

\begin{verbatim}
templateHTML: `<div><span>A</span></div>`
slots: { a: { kind: 'text', path: [0, 0] } }
\end{verbatim}

The hydration algorithm is:

\begin{enumerate}
  \item Read \texttt{serverRoot = host.firstElementChild}.
  \item Validate root tagName against the template (cheap check).
  \item Resolve every slot path against \texttt{serverRoot}.
  \item Validate slot node types, and element tagNames against the template DOM.
  \item Return a \texttt{HydratedBlock} whose \texttt{update()} writes into those resolved nodes.
\end{enumerate}

Two details in \texttt{runtime/hydration.ts} make this efficient:

\begin{itemize}
  \item Root tag validation is a fast regex over \texttt{templateHTML} (\texttt{getRootTagName}). This allows an early bail-out without parsing a template DOM.
  \item Slot resolution is just a \texttt{childNodes} index walk (\texttt{resolvePath}). It’s similar to \texttt{runtime/block.ts}, but hydration keeps its own copy with different error strings.
\end{itemize}

Spec: \texttt{hydrates without rewriting DOM when markup matches}.

What the spec is really asserting is \textbf{node identity preservation}. You can think of hydration as creating a \emph{mapping} from (slot key \(\rightarrow\) DOM node reference) without changing the structure:

\begin{itemize}
  \item before hydration: DOM nodes already exist (from server HTML)
  \item after hydration: we have slot node references pointing at those exact nodes
  \item on update: we mutate those nodes in place
\end{itemize}

\begin{center}
\begin{tikzpicture}[
  node distance=8mm,
  box/.style={draw, rounded corners, align=left, inner sep=5pt},
  arrow/.style={-Latex, thick}
]
\node[box] (server) {\textbf{Server DOM}\\existing nodes in host};
\node[box, right=18mm of server] (refs) {\textbf{Hydration}\\resolvePath(slot.path)\\slotNodes map};
\node[box, right=18mm of refs] (update) {\textbf{update(values)}\\mutate in place\\no replace};

\draw[arrow] (server) -- node[midway, above] {no rewrite} (refs);
\draw[arrow] (refs) -- (update);
\end{tikzpicture}
\end{center}

That test captures node identities before hydration and asserts identity stability after calling \texttt{hydrateBlock}:

\begin{itemize}
  \item same root element
  \item same inner \texttt{<span>}
  \item same text node
\end{itemize}

Then it updates slot values and checks the content changed without replacing nodes.

\section{Mismatch detection: when to bail out}

Hydration needs a mismatch policy. In v1 it is intentionally conservative.

\subsection{Root tag mismatch}

The first check is easy: compare the root tagName.

Implementation: \texttt{getRootTagName(templateHTML)} uses a small regex for the first tag.
If root tags don’t match, hydration does:

\begin{verbatim}
host.innerHTML = ''
return mountBlock(def, host, initialValues)
\end{verbatim}

This is the first key safety property: mismatched DOM never becomes a "half-hydrated" zombie tree. We either attach cleanly or we tear down and replace.

This strict policy gives you an important debugging invariant:

\begin{quote}
If hydration succeeds, later slot writes will never throw because of a missing node or wrong element type (for the validated paths).
\end{quote}

It’s also why HRBR v1 hydration doesn’t try to be clever about whitespace or comment-node drift: the moment your markup generation stops matching the compiler’s structural expectations, hydration deliberately stops and remounts.

Spec: \texttt{mismatch triggers remount for subtree}.

\subsection{Slot path mismatch: element tagName mismatch}

Root tag matching is not enough. You can have the same root element but different internal structure.

The conservative check in \texttt{runtime/hydration.ts} is:

\begin{itemize}
  \item Parse the expected template root into a detached DOM node.
  \item For each non-text slot:
    \begin{itemize}
      \item resolve the same path against the expected template DOM
      \item assert both are Elements
      \item assert \texttt{expected.tagName === actual.tagName}
    \end{itemize}
\end{itemize}

Why compare tag names at slot paths (instead of deep-walking the entire tree)?

\begin{itemize}
  \item Slots are where we will write later, so validating those boundaries catches many mismatches that would otherwise crash later.
  \item It avoids paying for a full tree diff during hydration.
  \item It keeps the rule extremely easy to explain and test.
\end{itemize}

If any assertion fails, hydration catches and bails out with client remount.

Spec: \texttt{slot path element mismatch triggers bailout remount (conservative)}.

\section{Why this design? (tradeoffs, alternatives, and invariants)}

Hydration is an area where it’s easy to ship something that works on your demo and breaks in production.
So I made a deliberate v1 choice: \textbf{strict validation + bailout}.

Here are the tradeoffs I’m optimizing for:

\begin{itemize}
  \item \textbf{Invariant: “hydrated means safe to update.”}
  If hydration returns successfully, the slots you’ll write to later have been validated against the template.

  \item \textbf{Low implementation risk.}
  Partial hydration systems can be correct, but they require more state and more surface area for mismatches.

  \item \textbf{Predictable failure mode.}
  When something is off (wrong tag, wrong path, wrong slot kind), we don’t keep walking a tightrope.
  We clear and remount.
\end{itemize}

Alternatives (all valid, just not v1 Relax):

\begin{itemize}
  \item \textbf{Deep tree diff during hydrate:} robust but expensive and complex.
  \item \textbf{Marker-based hydration:} embed explicit markers in HTML (comment anchors / data attributes). This is a common approach, and it’s likely the direction for a more advanced Relax hydration.
  \item \textbf{Island / partial hydration:} hydrate only interactive regions. Great UX story, more tooling.
\end{itemize}

For a framework-agnostic overview of hydration strategies and tradeoffs, see \href{https://web.dev/rendering-on-the-web/}{https://web.dev/rendering-on-the-web/}.

The test intentionally renders:\\
\texttt{<div><em>X</em></div>}\\
but the template expects:\\
\texttt{<div><span></span></div>}\\
So hydration remounts and you end up with a \texttt{<span>}.

\section{Hydrated update semantics: same slot kinds as mountBlock}

Hydration returns an object that looks like a mounted block:

\begin{itemize}
  \item \texttt{root} is the server root element
  \item \texttt{slotNodes} maps slot keys to resolved DOM nodes
  \item \texttt{update(values)} writes into those nodes
  \item \texttt{destroy()} removes hydrated listeners and removes \texttt{root}
\end{itemize}

The update logic supports the full set of slot kinds:

\begin{itemize}
  \item \texttt{text}: write \texttt{Text.nodeValue}
  \item \texttt{attr}: boolean attribute semantics and SVG xlink namespace
  \item \texttt{prop}: input correctness for \texttt{value} / \texttt{checked}
  \item \texttt{class}: string or array join
  \item \texttt{style}: string or object diff
  \item \texttt{event}: bind event handler(s)
\end{itemize}

Spec: \texttt{hydrates with full slot semantics (event/attr/prop/SVG) without rewriting existing DOM}.

Most of the slot patching logic is a copy of the mount-time block patching rules (Chapter 14), but hydration differs in one important place: event slots.

That test verifies:

\begin{itemize}
  \item DOM identity is preserved for root/button/input/link
  \item clicking calls the hydrated handler
  \item boolean attribute \texttt{hidden} is present
  \item \texttt{checked} property toggles
  \item SVG \texttt{xlink:href} is set in the xlink namespace
\end{itemize}

\subsection{Redundant write skipping}

Hydration tracks \texttt{prevValues} and skips updates when \texttt{Object.is(prev, next)}.
This matches the block runtime’s "skip redundant writes" optimization.

This is a subtle performance win because it also prevents unnecessary listener churn (given hydration’s v1 event strategy below).

\section{Event slots: v1 difference vs mountBlock}

There’s an intentional simplification in hydration:

\begin{itemize}
  \item \texttt{mountBlock} binds once and swaps \texttt{rec.current} (no listener churn).
  \item \texttt{hydrateBlock} removes the previous listener and re-adds on update.
\end{itemize}

You can see this directly in \texttt{runtime/hydration.ts}:

\begin{itemize}
  \item On \texttt{event} updates, it looks up \texttt{listenersByKey[k]}.
  \item If present, it calls \texttt{removeEventListener(prev.type, prev.handler)}.
  \item If the new value is a function, it creates a fresh wrapper handler and calls \texttt{addEventListener(slot.name, handler)}.
\end{itemize}

This is not as optimal as the bind-once approach in \texttt{mountBlock}, but it is still correct and easy to reason about.

Two important correctness notes (easy to miss):

\begin{itemize}
  \item \textbf{Redundant writes matter more with this strategy.} If you didn’t skip redundant writes, a repeated \texttt{update()} could churn listeners (remove+add) even when the handler didn’t change.
  \item \textbf{Destroy must remove hydrated listeners.} Otherwise the host subtree could be detached but still reachable through listener closures.
\end{itemize}

\begin{center}
\begin{tikzpicture}[
  node distance=8mm,
  box/.style={draw, rounded corners, align=left, inner sep=5pt},
  arrow/.style={-Latex, thick}
]
\node[box] (mb) {\textbf{mountBlock event slot}\\bind once\\swap \texttt{rec.current}};
\node[box, below=of mb] (hb) {\textbf{hydrateBlock event slot}\\remove old listener\\add new wrapper};
\node[box, right=18mm of mb] (pro) {\textbf{pro}\\simple to reason about\\low state};
\node[box, right=18mm of hb] (con) {\textbf{con}\\more listener churn\\(still correct)};

\draw[arrow] (mb) -- (hb);
\draw[arrow] (mb) -- (pro);
\draw[arrow] (hb) -- (con);
\end{tikzpicture}
\end{center}

The key safety property is lifecycle cleanup: \texttt{destroy()} removes any listeners hydration attached, so even in the “remove + re-add” strategy you don’t leak handlers across updates.

It still avoids leaks because \texttt{destroy()} removes any listeners it attached.

This choice keeps hydration simpler while still correct. If hydration becomes a performance hot spot, it can be refactored to reuse the bind-once strategy.

\section{SSR helper: renderBlockToString()}

\texttt{runtime/ssr.ts} provides a minimal string renderer used mainly for tests.

Contract:

\begin{itemize}
  \item create a detached host element
  \item \texttt{mountBlock(def, host, initialValues)}
  \item read \texttt{host.innerHTML}
  \item destroy the mounted block
\end{itemize}

This is explicitly a v1 teaching tool:

\begin{itemize}
  \item It uses the browser’s DOM implementation to do serialization.
  \item It’s not the kind of streaming SSR engine you’d use in production.
\end{itemize}

But it has an important property: it guarantees that the HTML string shape matches \texttt{mountBlock} semantics, because it’s literally produced by the same code.

Spec: \texttt{renders a block to HTML string, hydrates it, and updates via mountCompiledBlock}.

This gives a full lifecycle demo:

\begin{enumerate}
  \item render HTML string from a block
  \item insert it into the document
  \item hydrate it without rewriting
  \item mount a compiled block on top of the hydrated DOM and update via signals
\end{enumerate}

\section{Takeaways}

\begin{itemize}
  \item Hydration v1 is conservative: it validates structure and bails out on mismatch.
  \item When markup matches, hydration preserves DOM identity (no rewrites).
  \item Hydrated updates support the full slot matrix (text/attr/prop/class/style/event + SVG xlink).
  \item SSR v1 uses the DOM to produce strings for tests; it’s simple and correct, not optimized.
\end{itemize}

\section{Walking the tests (the executable spec)}

Hydration and SSR are notoriously easy to get "mostly right" and still ship broken edge cases.
Relax locks down the v1 behavior with two test files.

\subsection{\texttt{runtime/\_\_tests\_\_/hydration.test.ts}}

This suite defines the hydration safety model.

\begin{itemize}
  \item \textbf{hydrates without rewriting DOM when markup matches}
  captures \texttt{root/span/text} node identities before hydration, hydrates with an initial value, and asserts:
  \begin{itemize}
    \item the same elements/text nodes are still present (identity preserved)
    \item only the content changes
  \end{itemize}

  \item \textbf{hydrates with full slot semantics (event/attr/prop/SVG) without rewriting existing DOM}
  verifies the “slot matrix” works even when DOM is pre-existing:
  \begin{itemize}
    \item event handler is attached and called
    \item boolean-ish attr (\texttt{hidden}) is present
    \item \texttt{checked} is set as a property
    \item SVG \texttt{xlink:href} is set in the xlink namespace
  \end{itemize}

  \item \textbf{mismatch triggers remount for subtree}
  uses the wrong root tag (\texttt{<section>}) and asserts hydration replaces the entire subtree.

  \item \textbf{slot path element mismatch triggers bailout remount (conservative)}
  uses the same root tag but the wrong inner tag at a slot path (\texttt{<em>} vs expected \texttt{<span>}).
  The test asserts that hydration bails out and remounts so the template structure is restored.
\end{itemize}

\subsection{\texttt{runtime/\_\_tests\_\_/ssr.test.ts}}

This test is a full lifecycle demo "in miniature":

\begin{enumerate}
  \item \textbf{render} a block to an HTML string via \texttt{renderBlockToString}
  \item insert the HTML into a host
  \item \textbf{hydrate} without changing root/text node identities
  \item \textbf{mountCompiledBlock} to wire signals and update the hydrated DOM
\end{enumerate}

It’s intentionally pragmatic: rather than relying on timing, it calls \texttt{mounted.update({ value: 'B' })} to assert the DOM is updated.
The key idea being demonstrated is: \emph{hydration produces a block-shaped surface that the rest of HRBR can build on}.

\section{Online resources}

\begin{itemize}
  \item \href{https://developer.mozilla.org/en-US/docs/Web/HTML/Element/template}{MDN: \texttt{<template>}} (hydration compares template structure by parsing)
  \item \href{https://developer.mozilla.org/en-US/docs/Web/API/Element/firstElementChild}{MDN: \texttt{firstElementChild}} (the attach point)
  \item \href{https://developer.mozilla.org/en-US/docs/Web/API/Node/childNodes}{MDN: \texttt{Node.childNodes}} (path-based DOM addressing)
  \item \href{https://developer.mozilla.org/en-US/docs/Web/API/EventTarget/addEventListener}{MDN: events + listener removal}
  \item \href{https://developer.mozilla.org/en-US/docs/Web/SVG/Attribute/xlink:href}{MDN: \texttt{xlink:href}} (namespaced SVG attribute semantics)
  \item A broader hydration overview (framework-agnostic):
  \href{https://web.dev/rendering-on-the-web/}{web.dev: Rendering on the Web}.
\end{itemize}

\section{Common pitfalls}

These are the mistakes I see most often when people attempt hydration systems (or when I debug my own early versions):

\begin{itemize}
  \item \textbf{Treating whitespace as "not real DOM".} Paths are over \texttt{childNodes}, so whitespace text nodes count.
  \item \textbf{Hydrating when the template changed.} If \texttt{templateHTML} is an ABI, you must deploy server+client together.
  \item \textbf{Forgetting idempotence.} Hydration is often retried; cleanup must be safe even if called twice.
  \item \textbf{Listener leaks.} If you attach listeners during hydration, \texttt{destroy()} must remove them.
\end{itemize}

\section{Mini-exercise}

Open \texttt{runtime/hydration.ts} and make the mismatch error messages more actionable.

\begin{itemize}
  \item Add the slot key and the path (like \texttt{[2,0]}) in the "invalid path" assertion.
  \item Ensure the tests in \texttt{runtime/\_\_tests\_\_/hydration.test.ts} would still pass conceptually.
\end{itemize}

The goal is to practice the key engineering habit for hydration: when things go wrong, logs must tell you \emph{exactly} what structural assumption failed.
